\scanpage{462}\mspage{374}
% archive sheet number: 461
Donc on trouve
\[
\underline{\mathrm{Hom}}_{!\pi_0}(M,N)\xrightarrow{\approx}\underline{\mathrm{Hom}}_{\Phi}(C,N),
\]
et le corollaire 4.14.2. \unsure{donne} la conclusion voulue.

\uline{Corollaire} 4.14.4. Soient $M,N$ deux catégories accessibles, $\Phi$ et $\Psi$ deux ens. de catégories comme dans \unsure{4.14.1|4.14.2}, \deleted{\unsure{$N$} \illegible{1} \unsure{catégorie}} \inserted{on suppose que} \inserted{$M$ et} \deleted{\unsure{$C$}} $N$ \deleted{est} \inserted{satisfont $\unsure{\mathrm L_{\pi_0}|\mathrm L\pi_0}$, et \deleted{\unsure{sont}}} stables par $\varinjlim$ de type $\Phi$ et par $\varprojlim$ de type $\Psi$, Soit $\underline{\mathrm{Hom}}_{\Phi,\pi_0}^{\Psi}(M,N)$ la sous-catégorie str. pl. de $\underline{\mathrm{Hom}}_{\deleted{\unsure{\pi}}}(M,N)$ \unsure{formée|formé} des foncteurs $\pi_0$-accessibles, \deleted{et tels que} qui commutent aux $\varinjlim$ de type $\Phi$, et aux $\varprojlim$ de type $\Psi$. Alors $\underline{\mathrm{Hom}}_{\Phi,\pi_0}^{\Psi}(M,N)$ est accessible.

\uline{Dém} : Cette fois, on choisit $\pi>\pi_0$ tel qu'on ait $\operatorname{card}\unsure{I_{\mathrm{red}}^{.}|I\mathrm{red}^{.}|I_{\mathrm{red}}|I\mathrm{red}}\leq\unsure{\pi}$ pour tt $I$ dans $\Phi$ ou dans $\Psi$. La catégorie envisagée s'interprète alors (posant encore $C=M_{\pi}$) comme \deleted{\unsure{$K_{\alpha}$|la}}
